\chapter{Guided Solutions}
\label{ch:guided-solutions}

\begin{plainlanguage}
These solutions model reasoning, not merely final answers. Try each exercise
before reading its response. More than one wording or query can be correct when
its grain, assumptions, and validation are explicit.
\end{plainlanguage}

\section{Course orientation}

\subsection{Solution to the orientation exercise}
\label{sol:orientation}

A defensible workflow begins with the question and source documentation because
those establish what a row can mean. It then audits the source files, imports
them into a raw layer without silent reinterpretation, and converts them into
typed staging and core relations. Quality checks test row counts, keys,
relationships, ranges, and missingness before analysis. SQL results are
reconciled against known totals and described at their actual grain. Only then
should a finding be communicated, together with its limits. This sequence is
transferable because it relies on evidence and relational structure rather than
on OULAD-specific commands.

\section{Learning analytics and the problem}

\subsection{Solution to the classification check}
\label{sol:learning-analytics}

``Withdrawn attempts account for 31.2\% of recorded attempts'' is descriptive.
``Students with low activity are more likely to withdraw'' is predictive or
associational unless a causal design is supplied. ``Sending a message reduces
withdrawal'' is causal and requires a comparison strategy that addresses what
would have happened without the message. The important move is to classify the
claim before choosing methods.

\subsection{Solution to the question-design exercise}
\label{sol:learning-question}

One bounded question is: ``Among recorded student--module presentation
attempts, how does the distribution of final results vary by module and
presentation?'' The unit is one attempt. The population is the attempts present
in OULAD, not all distance-learning students. The outcome is
\texttt{final\_result}. A later analysis must consider whether presentations
differ in composition and timing. This question supports description; it does
not identify why outcomes differ.

\section{The OULAD data}

\subsection{Solution to the source-grain map}
\label{sol:source-grains}

\small
\begin{longtable}{@{}p{0.23\textwidth}p{0.42\textwidth}p{0.28\textwidth}@{}}
\caption{Suggested grains and candidate keys}\label{tab:solution-grains}\\
\toprule
Source & One row represents & Candidate key \\
\midrule
\endfirsthead
\toprule
Source & One row represents & Candidate key \\
\midrule
\endhead
\projectfile{courses.csv} &
One module presentation &
\texttt{module + presentation} \\
\projectfile{assessments.csv} &
One assessment attached to a module presentation &
\texttt{assessment\_id} \\
\projectfile{vle.csv} &
One VLE activity type within a module presentation &
\texttt{site\_id + module + presentation} \\
\projectfile{studentInfo.csv} &
One student--module presentation attempt &
\texttt{student\_id + module + presentation} \\
\projectfile{studentRegistration.csv} &
One registration record for an attempt &
\texttt{student\_id + module + presentation} \\
\projectfile{studentAssessment.csv} &
One student's submission record for one assessment &
\texttt{assessment\_id + student\_id} \\
\projectfile{studentVle.csv} &
One student's summarized interaction with one VLE site on one relative day &
\texttt{site\_id + student\_id + module + presentation + date} \\
\bottomrule
\end{longtable}
\normalsize

These are analytical candidate keys and should be tested, not trusted merely
because they sound plausible.

\subsection{Solution to the attempt trace}
\label{sol:attempt-trace}

Begin with one attempt in \projectfile{studentInfo.csv}. Match its registration
record using student, module, and presentation. Find the presentation's
assessments, then connect the student's submissions through
\texttt{assessment\_id} and \texttt{student\_id}. Separately, match VLE activity
through student, module, and presentation, with site and relative day adding
detail. Keep assessment and VLE facts in separate aggregates until each has
been reduced to attempt grain. Joining both detailed tables directly can
multiply rows.

\section{Relational foundations}

\subsection{Solution to the dangerous join}
\label{sol:dangerous-join}

Suppose an attempt has four assessment submissions and 300 daily-site activity
rows. A direct join between both detail tables can produce 1,200 rows for that
attempt. Summing clicks after that join would repeat each activity row four
times. The repair is to aggregate assessments to one row per attempt, aggregate
activity to one row per attempt, validate both intermediate keys, and only then
join the two attempt-level relations.

\subsection{Solution to the database-rules exercise}
\label{sol:database-rules}

A useful rule set includes:
\begin{itemize}
  \item a primary or unique key at the declared grain;
  \item foreign-key checks where the project intentionally enforces them;
  \item nonnegative click counts;
  \item accepted final-result values;
  \item explicit nullability for fields such as assessment due date; and
  \item row-count and orphan checks after every material transformation.
\end{itemize}
Database constraints prevent some invalid states. Quality queries remain
necessary for conditions that are contextual, cross-table, or deliberately
profiled instead of rejected.

\section{PostgreSQL setup}

\subsection{Solution to the environment exercise}
\label{sol:environment}

The environment record should identify Docker or local PostgreSQL, show the
client and server versions, name the database, and include a successful
\texttt{current\_database()} query. For Docker, the Compose configuration is
the authority for service, role, password, and database values. For a trusted
local installation, \texttt{postgresql:///learning\_analytics} can use the
current operating-system user. The repository's default URL names a
\texttt{learning\_analytics} role, while the local database target may create
only the database. That distinction should be diagnosed, not papered over.

\section{Raw import}

\subsection{Solution to the reconciliation exercise}
\label{sol:reconcile}

Create a three-column comparison for audited CSV, raw table, and typed core
relation. For the source used in this project, all seven CSV counts should be
accounted for. A mismatch is a stop condition until explained. First preserve
the import log. Then confirm that the audited file and delimiter are the ones
actually loaded. Finally, rerun only the smallest failing stage with
stop-on-error enabled and inspect rejected or malformed records. Never adjust a
target count simply to match an incomplete load.

\subsection{Solution to the failed-cast exercise}
\label{sol:failed-cast}

Keep the raw text row unchanged. Query the candidate column with a narrow
pattern or conversion guard to identify values that cannot be cast. Compare
them with the data dictionary and neighboring records. Then decide whether the
value represents a documented missing marker, a source anomaly, or a parsing
error. Implement the rule in staging, add a quality test, and rerun from a
reproducible boundary. Editing the raw table would destroy evidence of the
source condition.

\section{Data validation}

\subsection{Solution to the validation specification}
\label{sol:validation-spec}

For \texttt{student\_vle}, a compact specification includes:
\begin{itemize}
  \item row count reconciles to the audited source;
  \item the declared composite key is unique;
  \item every module-presentation pair exists in courses;
  \item every site belongs to the corresponding VLE presentation;
  \item every student attempt exists at attempt grain;
  \item \texttt{sum\_click} is nonnegative; and
  \item activity-day range and pre-start activity are profiled.
\end{itemize}
Pre-start records are a warning profile because they may be legitimate access,
not automatically invalid rows.

\subsection{Solution to the orphan-debugging exercise}
\label{sol:orphan-debug}

Start with an anti-join that returns the orphaned child keys and a small sample,
not only a count. Confirm that the join uses the full composite relationship;
omitting presentation is a common mistake. Compare data types and whitespace
normalization. Then inspect whether the parent source lacks the key or whether
staging transformed it incorrectly. Fix the earliest reproducible cause and
retain the anti-join as a regression test.

\section{Introductory SQL}

\subsection{Solution to the inspection exercise}
\label{sol:sql-inspection}

\begin{lstlisting}[language=SQL]
SELECT
    student_id,
    module,
    presentation,
    final_result
FROM core.student_info
WHERE final_result IN ('Fail', 'Withdrawn')
ORDER BY module, presentation, student_id
LIMIT 25;
\end{lstlisting}

This is an inspection query. The limit makes the display manageable but does
not define the analytical population.

\subsection{Solution to the outcome exercise}
\label{sol:sql-outcomes}

\begin{lstlisting}[language=SQL]
SELECT
    module,
    final_result,
    COUNT(*) AS attempts
FROM core.student_info
GROUP BY module, final_result
HAVING COUNT(*) >= 100
ORDER BY module, attempts DESC;
\end{lstlisting}

The result is at module-by-outcome grain. The \texttt{HAVING} clause filters
groups after aggregation; a \texttt{WHERE COUNT(*)} expression would be
invalid.

\subsection{Solution to the missingness exercise}
\label{sol:sql-missing}

\begin{lstlisting}[language=SQL]
SELECT
    COUNT(*) AS attempts,
    COUNT(imd_band) AS recorded_imd_band,
    COUNT(*) FILTER (WHERE imd_band IS NULL) AS missing_imd_band,
    ROUND(
        100.0 * COUNT(*) FILTER (WHERE imd_band IS NULL) / COUNT(*),
        1
    ) AS missing_pct
FROM core.student_info;
\end{lstlisting}

\texttt{COUNT(imd\_band)} ignores null values, while \texttt{COUNT(*)} counts
rows. The expected missing count is 1,111 attempt records.

\subsection{Solution to the repeat-appearance exercise}
\label{sol:sql-repeats}

\begin{lstlisting}[language=SQL]
SELECT
    student_id,
    COUNT(*) AS attempt_count,
    COUNT(DISTINCT module) AS module_count
FROM core.student_info
GROUP BY student_id
HAVING COUNT(*) > 1
ORDER BY attempt_count DESC, student_id
LIMIT 25;
\end{lstlisting}

This query identifies people with multiple attempt rows. It does not by itself
tell whether the rows are retries, different modules, or another pattern; that
requires examining module and presentation.

\section{Cumulative lab}

\subsection{Solution to lab checkpoint 1}
\label{sol:lab-grain}

\texttt{COUNT(*)} counts the 32,593 rows in the attempt-level relation, so its
unit is attempts. \texttt{COUNT(DISTINCT student\_id)} counts the 28,785 unique
student identifiers represented by those rows. The attempt count is larger
because some students appear in more than one module presentation. Neither
measure is ``more correct'' without a question that names its unit.

\subsection{Solution to lab checkpoint 2}
\label{sol:lab-environment}

For the Docker path, confirm the Compose service is running and use the role,
password, port, and database declared there. Recreating the service from the
project configuration may be appropriate when no persistent work would be
lost. For a trusted local path, use the current local role with
\texttt{postgresql:///learning\_analytics}, or deliberately create and grant
the named role. Import scripts cannot fix the error because authentication
happens before SQL or CSV loading begins.

\subsection{Solution to lab checkpoint 3}
\label{sol:lab-reconcile}

First save the failed command, full log, audited file checksum, and observed
counts. Second, use a clean raw table and rerun the single
\texttt{studentAssessment.csv} load with stop-on-error behavior, watching for
delimiter, quoting, encoding, or newline errors. Third, compare the final file
record with the loaded tail and verify that the source audit counted data rows
the same way as PostgreSQL. Repair the loader or source acquisition step; do
not invent the missing row.

\subsection{Solution to lab checkpoint 4}
\label{sol:lab-quality}

One acceptable note is:
\begin{quote}
All 18 error-level tests passed, so the current build clears the project's
blocking quality gate. The suite also profiles 688,988 VLE activity rows before
presentation day zero; these are retained, but later feature windows must state
whether they include them. Eleven assessments have no due date, so deadline-
relative features need a missingness rule. IMD band is missing for 1,111
attempts, so demographic summaries must report their denominator and missing
category. These warnings document consequential source conditions rather than
failed integrity tests.
\end{quote}

\subsection{Solution to lab checkpoint 5}
\label{sol:lab-synthesis}

\begin{quote}
The data contain 32,593 student--module presentation attempts contributed by
28,785 distinct students, so several students appear more than once. Pass is
the most common recorded result (12,361 attempts), while 10,156 attempts
(31.2\%) are recorded as Withdrawn. These values describe attempts, not unique
people, and should not be interpreted as causal explanations of learning
outcomes. The blocking quality checks pass, but 688,988 activity rows occur
before presentation day zero; any later engagement window must decide and
document whether those rows are included. Assessment due-date and IMD-band
missingness also constrain deadline-based or demographic comparisons. A useful
next question is whether the outcome distribution differs across module
presentations after first checking that each comparison uses consistent
attempt-level denominators.
\end{quote}

\section{How to use these solutions}

Compare the structure of your reasoning with the examples. A strong response
usually names the unit, the population, the operation, the check, and the
limit. Rewording is encouraged; unexamined copying is not.
